% =====================================================================
%  Section for the main paper: T = p = 4 application to the ENIA census.
%  Requires: booktabs, amsmath, graphicx.
%  Figures: fig_scatter.pdf, fig_ghat.pdf
% =====================================================================

\section{Application: heterogeneous input elasticities in Chilean
manufacturing}
\label{sec:chile}

\subsection{Specification}

Production function estimation is the canonical setting in which a coefficient
that is almost certainly heterogeneous is nevertheless restricted to be common
across units. In the control-function approach of \citet{OlleyPakes1996} and
\citet{LevinsohnPetrin2003}, and in the refinements of \citet{ACF2015} and
\citet{GNR2020}, the entire apparatus of proxy variables is devoted not to the
elasticities but to the \emph{intercept}: the unobserved productivity term
$\omega_{it}$ that is correlated with input choices and generates the
transmission bias. The elasticities themselves are assumed common.

The correlated random coefficient model inverts this division of labour. We
estimate the gross-output production function
\begin{equation}
  y_{it} \;=\; X_{it}'(\beta_i + \delta_t) \;+\; \varepsilon_{it},
  \qquad
  X_{it} \;=\; (1,\; k_{it},\; l_{it},\; x_{it})',
  \qquad t = 1,\dots,4,
  \label{eq:klm}
\end{equation}
where $y$ is log real gross output, $k$ and $l$ are logs of the real capital
stock and employment, and $x$ is the log of real intermediate inputs,
constructed as the sum of deflated materials and energy. This is the
single-intermediate-input form used by \citet{GNR2020}; aggregating materials
and energy is not a concession here but a choice we return to in
Section~\ref{sec:designsize}, since it materially improves what the design can
resolve.

The plant-specific intercept $\beta_{i1}$ absorbs $\omega_i$ directly, so no
proxy variable and no monotonicity condition is required. In exchange, all
three elasticities may vary across plants and be correlated with the input
history. The vector $\delta_t$ is a common drift that shifts the entire
coefficient vector in period $t$, allowing the technology itself to move over
the window while $\beta_i$ carries the permanent plant-specific component. We
normalise $\delta_1 = 0$, so the estimand $\mathbb{E}[\beta_i]$ is the average
elasticity vector anchored at the first year of the window.

With $p = 4$ coefficients the model is just identified at $T = 4$, which places
it in the irregular case of Section~3. Writing $X_i$ for the $4 \times 4$
design matrix whose rows are $X_{it}'$, and $W_i$ for the $T \times Q$ matrix
whose $t$th row places $X_{it}'$ in block $t-1$ and is zero elsewhere
($Q = (T-1)p = 12$), the determinant
\begin{equation*}
  h_i \;=\; |\det(X_i)|
\end{equation*}
governs identification. A plant with $h_i$ near zero has moved its input bundle
along a nearly degenerate path over the four years, and its coefficient vector
is not separately identified. The economic content of the model is that such
plants differ systematically from the rest: a plant reallocates inputs
precisely when its own returns to doing so are high, so
$g(h) = \mathbb{E}[\beta_i \mid h_i = h]$ is not constant and the trimmed
estimator of \citet{GrahamPowell2012} is biased for the average partial effect.

Estimation proceeds in two steps. Since
$\operatorname{adj}(X_i)X_i = \det(X_i)I$, premultiplying \eqref{eq:klm} by
$\operatorname{adj}(X_i)$ annihilates $\beta_i$ whenever $\det(X_i) = 0$,
leaving
\begin{equation}
  \operatorname{adj}(X_i)\,Y_i \;=\; \operatorname{adj}(X_i)\,W_i\,\delta
                                 \;+\; \operatorname{adj}(X_i)\,\varepsilon_i .
  \label{eq:driftstep}
\end{equation}
The drift is therefore identified off the stayers alone, and $\hat\delta$ is
the least squares coefficient in \eqref{eq:driftstep} over units with
$h_i \leq h_0$. The coefficient step then runs on the drift-purged outcome
$\tilde Y_i = Y_i - W_i\hat\delta$, giving
$\hat\beta_i = \operatorname{adj}(X_i)\tilde Y_i/\det(X_i)$ for movers.

The trimming threshold thus does double duty: it defines the mover sample and
simultaneously the stayer sample identifying $\delta$. Lowering $p_0$ reduces
the bias of the coefficient step but shrinks the sample available for the drift
step, and since $\hat\delta$ enters $\tilde Y_i$ for \emph{every} plant, noise
in the drift step contaminates movers as well as stayers. This tension is the
organising constraint of the section, and it is why the size of the design is
not innocuous: $Q$ grows as $(T-1)p$, so each additional regressor in a
just-identified design adds $2p-1$ parameters that must be recovered from a
shrinking subsample.

\subsection{Data and design}

We use the full Encuesta Nacional Industrial Anual, the annual census of
Chilean manufacturing plants with at least ten employees, covering 1979--1996:
86{,}186 plant-year observations on 10{,}927 plants. Real gross output, real
materials, real energy and the real capital stock (buildings, machinery and
vehicles, mid-year) are sectorally deflated series; the capital stocks are
those of \citet{Greenstreet2007}, built by a permanent inventory procedure
separately for structures, equipment and transportation. Labour is the total
employment headcount. We access the census through the replication package of
\citet{Raval2023}, and depart from his construction in using the mid-year real
capital stock rather than a rental-cost aggregate, since the auxiliary
depreciation and interest-rate files his code merges are not needed for our
purpose.

Estimating \eqref{eq:klm} requires plants observed in four \emph{consecutive}
years with strictly positive output and all inputs. Coverage is remarkably flat
across the sample period, with every window from 1987 onwards retaining between
3{,}250 and 3{,}414 plants; we use 1992--1995, which retains the most
(3{,}414) and sits in the post-stabilisation period rather than spanning the
debt crisis. Because each plant contributes a single observation, the units are
independent and no clustering is required. Table~\ref{tab:design} reports the
design diagnostics of Section~2.

\begin{table}[t]
\centering
\caption{Design diagnostics, ENIA 1992--1995}
\label{tab:design}
\begin{tabular}{lc}
\toprule
Plants ($N$), each observed $T = 4$ years        & 3{,}414 \\
Coefficients ($p$)                               & 4 \\
Drift parameters ($Q = (T-1)p$)                  & 12 \\
$\pi_0 = \Pr(\det X_i = 0)$ (exact atom)         & 0.0076 \\
Median $h_i$                                     & $5.97 \times 10^{-3}$ \\
$h_i$ at the 10th percentile                     & $5.05 \times 10^{-4}$ \\
Median $\operatorname{cond}(X_i)$                & $5.9 \times 10^{3}$ \\
$\hat a$ in $f(h) \sim h^{a}$                    & $-0.320$ \\
\bottomrule
\end{tabular}

\vspace{0.6em}
\begin{minipage}{0.92\textwidth}
\footnotesize \emph{Notes.} $\hat a$ is the slope of $\log \hat f(h)$ on
$\log h$ over the lowest three deciles of $h$, from a reflected Epanechnikov
density estimate.
\end{minipage}
\end{table}

Two features matter. First, $\hat a \approx 0$: the density of the determinant
is bounded away from zero at the boundary, the thick-boundary regime $a < 1$ of
Theorem~1 in which the trimmed estimator converges more slowly than $\sqrt N$.
Second, the point mass at $\det X_i = 0$ is negligible. Requiring a
$4 \times 4$ design to be exactly singular is far more demanding than requiring
one input to be unchanged between two adjacent years, and under $1\%$ of plants
satisfy it. Assumption~4(d) is therefore not in question, and the trimming
proportion is a free choice rather than something the data forces.

We use that freedom to keep it small. Every estimate in Table~\ref{tab:main}
trims at most $10\%$ of plants. The bias of the trimmed estimator is
proportional to $p_0$, so holding $p_0$ below one tenth keeps the estimand
close to the average partial effect over the whole population rather than over
an endogenously selected subpopulation of movers. The costs are the ones noted
above---retained ill-conditioned designs, and a thinner stayer sample for the
drift step---and we let the standard errors display them rather than buying
precision by trimming a quarter of the sample.

\subsection{Is there heterogeneity to speak of?}

Equation~\eqref{eq:klm} is exactly identified, so within-plant residuals vanish
and the model cannot test its own homogeneity restriction. We add a fifth year,
1996, which leaves 3{,}158 plants with one residual degree of freedom apiece.
Estimating $\delta$ by pooled least squares on the interacted design, purging
it from the outcome and then fitting plant-specific regressions gives
$\hat\sigma = 0.130$, and the \citet{Swamy1970} statistic for
$H_0: \beta_i = \beta$ for all $i$ is $S = 1.19 \times 10^{5}$ on $12{,}580$
degrees of freedom. Homogeneity is rejected overwhelmingly. The dispersion is
economically large as well: among movers the interquartile range of
$\hat\beta_i$ is $1.19$ for intermediates and $2.15$ for capital, and $0.93$
and $1.39$ respectively in the better-conditioned half of the sample where
estimation error contributes less.

Heterogeneity alone does not require the machinery of this paper; only
heterogeneity \emph{correlated} with the determinant does. We regress the
estimated individual coefficients on $h_i$ among movers,
\begin{equation}
  \hat\beta_{i,p} \;=\; \alpha_p + \kappa_p h_i + e_{ip},
  \label{eq:kappa}
\end{equation}
weighting by $h_i^2$, which is proportional to the inverse conditional variance
when the $\Sigma_v(h_i)/h_i^{2}$ term dominates. Unweighted least squares is
uninformative here, being dominated by near-boundary observations that carry
almost no information about $g$.

\begin{table}[t]
\centering
\caption{Selection gradient $\hat\kappa_p$ in \eqref{eq:kappa}}
\label{tab:kappa}
\begin{tabular}{lccc}
\toprule
& $\hat\kappa$ & $t$ & $p$ \\
\midrule
Capital       & $-0.039$ & $-0.85$ & $0.393$ \\
Labour        & $-0.581$ & $-7.03$ & $<0.001$ \\
Intermediates & $\phantom{-}0.515$ & $\phantom{-}4.29$ & $<0.001$ \\
\bottomrule
\end{tabular}

\vspace{0.6em}
\begin{minipage}{0.92\textwidth}
\footnotesize \emph{Notes.} Movers only, trimming at $10\%$
($N_{\text{mov}} = 3{,}072$, $342$ stayers). Weighted by $h_i^2$;
heteroskedasticity-robust standard errors.
\end{minipage}
\end{table}

Labour and intermediates show sharp gradients of opposite sign; capital does
not. Plants whose input bundle moves least have markedly \emph{higher}
estimated labour elasticities and lower intermediate elasticities. This is the
configuration in which the trimmed estimator is biased and an extrapolation
estimator is called for.

Figure~\ref{fig:scatter} plots the individual coefficient estimates against
$h_i$ with the trimmed region shaded, and Figure~\ref{fig:ghat} the local
linear estimate of $g$ with its extrapolation to the boundary. Both display the
$h_i^{-2}$ heteroskedasticity plainly: the cloud fans out sharply as $h_i$
approaches the threshold and the confidence band for $\hat g$ widens in the
same region. The labour panel of Figure~\ref{fig:ghat} shows the gradient of
Table~\ref{tab:kappa} directly, $\hat g$ rising steeply as $h \to h_0$.

\begin{figure}[t]
\centering
\includegraphics[width=\textwidth]{fig_scatter.pdf}
\caption{Individual coefficient estimates $\hat\beta_{i,p}$ against
$h_i = |\det X_i|$, ENIA 1992--1995. Shaded region: units trimmed at
$h_0 = q_{10}(h)$. Vertical axes truncated at the 97th percentile of
$|\hat\beta_{i,p}|$.}
\label{fig:scatter}
\end{figure}

\begin{figure}[t]
\centering
\includegraphics[width=\textwidth]{fig_ghat.pdf}
\caption{Local linear estimate of $g(h) = \mathbb{E}[\beta_i \mid h_i = h]$
among movers (solid), with 90\% bootstrap band, and its linear extrapolation
into the trimmed region (dashed). Shaded region as in
Figure~\ref{fig:scatter}.}
\label{fig:ghat}
\end{figure}

\subsection{The estimated drift}
\label{sec:drift}

Table~\ref{tab:drift} reports $\hat\delta_t$ at the $10\%$ threshold, where
$342$ stayers identify the twelve parameters.

\begin{table}[t]
\centering
\caption{Estimated coefficient drift $\hat\delta_t$, $10\%$ trimming
(1992 normalised to zero)}
\label{tab:drift}
\begin{tabular}{lcccc}
\toprule
& Intercept & Capital & Labour & Intermediates \\
\midrule
1993 & $-0.459$ & $-0.065$ & $\phantom{-}0.008$ & $0.091$ \\
     & (0.245) & (0.034) & (0.066) & (0.039) \\
1994 & $-1.270$ & $-0.108$ & $-0.082$ & $0.227$ \\
     & (0.429) & (0.052) & (0.121) & (0.065) \\
1995 & $-1.630$ & $-0.107$ & $-0.169$ & $0.286$ \\
     & (0.541) & (0.065) & (0.162) & (0.085) \\
\bottomrule
\end{tabular}

\vspace{0.6em}
\begin{minipage}{0.92\textwidth}
\footnotesize \emph{Notes.} $\hat\delta$ from \eqref{eq:driftstep}, estimated
on the $342$ units with $h_i \leq h_0 = q_{10}(h)$. Bootstrap standard errors
over plants, 400 replications.
\end{minipage}
\end{table}

The drift is real and economically legible. A Wald test of $\delta = 0$ against
the bootstrap covariance gives $23.94$ on twelve degrees of freedom
($p = 0.021$). The intermediate component rises monotonically---$0.091$,
$0.227$, $0.286$ across 1993--95, each more than two standard errors from
zero---while the capital and intercept components fall. Read literally, the
return to intermediate inputs rose over the window while the return to capital
declined, holding the plant-specific component fixed. Imposing
$\delta \equiv 0$ would therefore be a substantive restriction rather than a
simplification, which matters because the restriction is tempting: it is the
drift step, not the coefficient step, that makes low trimming expensive.

\subsection{Estimates}

\begin{table}[t]
\centering
\caption{Average input elasticities, ENIA 1992--1995}
\label{tab:main}
\begin{tabular}{lcccc}
\toprule
Trimmed & Stayers & Capital & Labour & Intermediates \\
\midrule
\multicolumn{5}{l}{\emph{Panel A: Graham--Powell trimmed estimator}} \\
\phantom{0}2\% & \phantom{0}69 & $0.175$ & $0.482$ & $0.047$ \\
               &     & (0.532) & (0.482) & (0.303) \\
\phantom{0}4\% & 137 & $0.297$ & $0.592$ & $0.294$ \\
               &     & (0.213) & (0.338) & (0.164) \\
\phantom{0}6\% & 205 & $0.205$ & $0.374$ & $0.334$ \\
               &     & (0.196) & (0.299) & (0.124) \\
\phantom{0}8\% & 274 & $0.107$ & $0.303$ & $0.286$ \\
               &     & (0.184) & (0.272) & (0.127) \\
10\%           & 342 & $0.145$ & $0.287$ & $0.380$ \\
               &     & (0.177) & (0.234) & (0.103) \\
\midrule
\multicolumn{5}{l}{\emph{Panel B: local polynomial correction}} \\
\phantom{0}2\% & \phantom{0}69 & $0.195$ & $0.527$ & $0.012$ \\
\phantom{0}4\% & 137 & $0.306$ & $0.561$ & $0.307$ \\
\phantom{0}6\% & 205 & $0.196$ & $0.391$ & $0.335$ \\
\phantom{0}8\% & 274 & $0.076$ & $0.358$ & $0.288$ \\
10\%           & 342 & $0.121$ & $0.354$ & $0.378$ \\
\midrule
\multicolumn{5}{l}{\emph{Panel C: common-coefficient benchmarks}} \\
Pooled OLS, 1992 base & & $0.112$ & $0.224$ & $0.715$ \\
Pooled OLS, year mean & & $0.112$ & $0.238$ & $0.704$ \\
Within, 1992 base     & & $0.044$ & $0.257$ & $0.486$ \\
\bottomrule
\end{tabular}

\vspace{0.6em}
\begin{minipage}{0.92\textwidth}
\footnotesize \emph{Notes.} $N = 3{,}414$ plants, $T = p = 4$. Local polynomial
estimator uses $m = 1$, $c = b/h_0 = 30$ and inverse-variance weights. Drift
estimated as in Table~\ref{tab:drift} and subtracted before the coefficient
step. Standard errors from 400 bootstrap replications over plants; Panel~B
estimates carry standard errors within $0.02$ of the corresponding Panel~A
entries and are omitted. Panel~C estimates
$y_{it} = X_{it}'(\beta + \delta_t)$ by pooled and within least squares. All
thresholds are at or below $10\%$ by design.
\end{minipage}
\end{table}

The intermediate-input elasticity is the coefficient this design resolves. From
$4\%$ trimming onwards it sits between $0.29$ and $0.38$ with standard errors
falling from $0.164$ to $0.103$, and it lies well below both the pooled OLS
value of $0.70$ and the within value of $0.49$. At the $10\%$ threshold the gap
to pooled OLS is more than three standard errors. Allowing the elasticity to be
plant-specific and correlated with input history therefore lowers the average
substantially relative to the common-coefficient estimate---which is the kind
of quantity the model exists to deliver, and which no amount of proxy-variable
machinery would recover, since the proxies are aimed at the intercept.

Capital and labour are estimated less precisely. Both remain positive and in
the neighbourhood of the benchmark values throughout, and neither shows the
sign instability that would signal a broken design, but with standard errors of
$0.18$ and $0.23$ at the $10\%$ threshold we would not report either as a point
estimate. The ordering is not accidental: the conditional variance of
$\hat\beta_i$ is governed by the corresponding row of
$\operatorname{adj}(X_i)$, and intermediate inputs are adjusted far more freely
across years than capital or headcount, so their coefficient is the one the
determinant can resolve.

The $2\%$ row should be read with care and we would not build on it. It trims
so little that only $69$ stayers remain to identify twelve drift parameters,
and its instability---the intermediate elasticity collapsing to $0.047$---is
a property of the drift step rather than of the coefficients.

Finally, under coefficient homogeneity the average partial effect and the
within estimator estimate the same object; under \emph{correlated}
heterogeneity they do not. Bootstrapping the difference at the $10\%$
threshold, no individual difference exceeds one standard error and the joint
Wald statistic is $1.64$ on three degrees of freedom ($p = 0.65$). Given the
standard errors in Table~\ref{tab:main} this is a statement about power rather
than about the absence of correlated heterogeneity: homogeneity is decisively
rejected, the gradients in Table~\ref{tab:kappa} are sharp for two of three
inputs, and the point estimate for intermediates differs from the within value
by $0.11$---but the design cannot convert that into a statistically
distinguishable revision of the average.

\subsection{Design size is a design choice}
\label{sec:designsize}

The specification in \eqref{eq:klm} aggregates materials and energy into a
single intermediate input. Because $Q = (T-1)p$, disaggregating them is far
from free: it raises the design from $T = p = 4$ to $T = p = 5$, the drift
parameters from twelve to twenty, and---since five consecutive years are then
required---lowers $N$ from 3{,}414 to 3{,}133 and the stayer count at
$p_0 = 0.10$ from $342$ to $314$. Estimating the larger model on the same
census and at the same thresholds, no elasticity is recovered with useful
precision: the materials coefficient moves between $0.15$ and $0.46$ across the
sub-$10\%$ grid with standard errors of $0.14$ to $0.39$, against $0.29$ to
$0.38$ with standard errors of $0.10$ to $0.16$ here.

This is worth stating as a practical lesson, because the instinct in applied
work runs the other way: a richer input list looks like a more faithful
technology. In a just-identified correlated random coefficient design it is not
free, and the price is paid twice---once in the conditioning of $X_i$ and again
in the number of drift parameters the stayers must support. Aggregation to a
composite that still \emph{moves} is preferable to disaggregation into
components that do not, which is also why dropping labour rather than
aggregating intermediates would be the wrong economy: labour is the input whose
stickiness generates near-singular designs in the first place.

That stickiness is visible throughout. Labour carries the sharpest selection
gradient of any coefficient in Table~\ref{tab:kappa} while having among the
least precise averages, and it is integer headcounts that generate what atom
remains at $\det X_i = 0$. In the two-period designs common in this literature,
where a single lumpy input is allowed to vary, the same stickiness produces
atoms above a quarter, $p_0$ is bounded below by $\pi_0$, and the trimmed
estimator is inconsistent rather than merely biased. The four-period census
design escapes that regime, and the diagnostics of Table~\ref{tab:design} are
what tell an applied researcher in advance which regime they are in.
